\documentclass[11pt,a4paper]{article}
\usepackage{xeCJK} 
\usepackage[T1]{fontenc}
\usepackage{geometry}
\geometry{left=2.5cm,right=2.5cm,top=2.5cm,bottom=2.5cm}
\usepackage{hyperref}
\usepackage{listings}
\usepackage{xcolor}
\usepackage{graphicx}
\usepackage{amsmath}
\usepackage{amssymb}
\usepackage{booktabs}
\usepackage{tikz}
\usetikzlibrary{shapes.geometric, arrows, positioning, shadows} % 修复：shadow -> shadows

\title{\textbf{ACE Agent (Automated Clustering Expert) \\ 生产级智能体聚类系统工程手册}}
\author{系统架构组 - Gemini CLI 实验室}
\date{2026年4月17日}

\definecolor{codegreen}{rgb}{0,0.6,0}
\definecolor{codegray}{rgb}{0.5,0.5,0.5}
\definecolor{codepurple}{rgb}{0.58,0,0.82}
\definecolor{backcolour}{rgb}{0.97,0.97,0.95}

\lstset{
    backgroundcolor=\color{backcolour},   
    commentstyle=\color{codegreen},
    keywordstyle=\color{blue},
    numberstyle=\tiny\color{codegray},
    stringstyle=\color{codepurple},
    basicstyle=\ttfamily\small,
    breakatwhitespace=false,         
    breaklines=true,                 
    captionpos=b,                    
    keepspaces=true,                 
    numbers=left,                    
    numbersep=5pt,                  
    showspaces=false,                
    showstringspaces=false,
    showtabs=false,                  
    tabsize=2,
    frame=single
}

\begin{document}

\maketitle

\begin{abstract}
本文档旨在为 ACE Agent 项目提供详尽的工程实现细节。ACE Agent 是一个基于 Orchestrator-Worker 架构的多智能体系统，专为自动化聚类分析设计。系统集成了语义路由、代码自愈、安全沙箱及学术报告生成等核心模块。通过本文档，开发者可以完整复现该系统，并理解大语言模型（LLM）如何在非确定性任务中实现工业级的稳定产出。
\end{abstract}

\tableofcontents
\newpage

\section{引言与背景}
传统的自动化机器学习（AutoML）工具通常依赖硬编码的流水线，难以处理模糊的用户意图。ACE Agent 的设计目标是利用 LLM 的逻辑推理能力，构建一个具备“闭环纠错”能力的自愈系统。

\section{系统整体架构 (Architecture Deep Dive)}

\subsection{Orchestrator-Worker 模式}
系统核心遵循编排器模式，分为四个逻辑平面：
\begin{enumerate}
    \item \textbf{意图解析面 (The Router)}: 将自然语言（NLP）映射为计算意图。
    \item \textbf{专家调度面 (The Orchestrator)}: 决定哪些专家参与任务，并管理记忆。
    \item \textbf{代码执行面 (The Sandbox)}: 提供物理隔离的执行环境。
    \item \textbf{成果综合面 (The Synthesizer)}: 聚合专家输出。
\end{enumerate}

\subsection{数据流拓扑图}
\begin{center}
\begin{tikzpicture}[node distance=2cm, auto]
    \tikzstyle{block} = [rectangle, draw, fill=blue!10, text width=6em, text centered, rounded corners, minimum height=3em, drop shadow]
    \tikzstyle{line} = [draw, -latex']
    
    \node [block] (user) {用户输入};
    \node [block, right of=user, node distance=4cm] (router) {MasterRouter};
    \node [block, below of=router] (orch) {Orchestrator};
    \node [block, left of=orch, node distance=4cm] (expert) {Expert Agents};
    \node [block, below of=expert] (sandbox) {CoderSandbox};
    \node [block, below of=orch] (report) {报告生成器};

    \path [line] (user) -- (router);
    \path [line] (router) -- (orch);
    \path [line] (orch) -- (expert);
    \path [line] (expert) -- (sandbox);
    \path [line] (sandbox) -- (expert);
    \path [line] (expert) -- (orch);
    \path [line] (orch) -- (report);
\end{tikzpicture}
\end{center}

\section{核心功能模块详解}

\subsection{语义路由与意图判定 (MasterRouter)}
\texttt{agent\_core/router.py} % 修复：添加转义下划线
模块不使用正则匹配，而是通过精心设计的 Prompt 引导 LLM 返回标准化的 JSON 对象。

\subsection{专家自愈循环 (Self-Correction Loop)}
在 \texttt{expert\_sub\_agents/base.py} 中实现的 \texttt{execute\_with\_self\_correction} 函数是系统的鲁棒性保障：
如果代码报错，系统将 Traceback 截取后反馈给专家 Agent。

\subsection{安全沙箱执行 (CoderSandbox)}
\texttt{tools/coder\_sandbox.py} 通过 \texttt{exec()} 函数执行代码。

\begin{lstlisting}[language=Python]
exec_env = {
    "X": dataset.X, 
    "y": dataset.y, 
    "np": np,       
    "plt": plt,
    "artifacts": {}, 
    "__builtins__": SAFE_BUILTINS 
}
\end{lstlisting}

\section{LLM 接入点与 Prompt 策略}

\subsection{MasterRouter 提示词工程}
核心在于“强制约束”：要求只返回 JSON 格式。

\subsection{专家代码生成提示词}
增加了“物理隔离约束”：禁止使用 \texttt{pd.read\_csv}。

\section{函数调用时序与协作}

\subsection{初始化阶段}
\texttt{web\_demo.py} 初始化 \texttt{UniversalLLMClient} 与 \texttt{ACESupervisor}。

\section{关键类与方法索引 (API Reference)}

\begin{itemize}
    \item \texttt{UniversalLLMClient.chat\_completion}: 统一封装了调用协议。
    \item \texttt{ACEJsonEncoder}: 解决 \texttt{numpy.ndarray} 序列化痛点。
    \item \texttt{LatexReportGenerator.generate}: 渲染专业的 \LaTeX 源码。
\end{itemize}

\section{复现与环境配置指南}

\subsection{环境配置}
复现时需注意抑制 Windows MKL 内存泄漏警告。

\subsection{中文字体支持}
在 Win32 环境下，必须显式设置中文字体。

\end{document}
